% ECONOMICS_PROBLEM: {"id":"C73-8","title":"Payoff functionals guaranteeing a value under one-sided information","jel_code":"C73","source_id":"OP-0167"}
% FIRST_POSTED: 2026-10-11
% ATTEMPT_STATUS: solved
% ENTRY_KIND: research_attempt
% SUBMISSION: {"kind":"research_attempt","category":"economics","problem_id":"C73-8","model":{"provider":"OpenAI","version":"gpt-6-astra","reasoning":"ultra (Codex runtime metadata)","generated_on":"2026-10-11"},"other_models":"OpenAI gpt-6-astra high-reasoning review tasks; a ChatGPT Pro response (underlying version and reasoning not exposed); Claude Opus 5.5 and Claude Fable 5.1 (Claude Code), credited by the pinned Lean companion for its development. Original v1 Codex model version not exposed.","tools":"Primary-source browsing; GitHub CLI; desktop LaTeX compiler; Tectonic; Poppler PDF rendering. No Lean build or verification of added theorems.","target":"For arbitrary fixed finite nonempty K,I,J, intrinsically characterize all bounded Borel payoffs admitting a value at every prior in the exact one-sided-information simultaneous-action game.","scope":"Theorems 2.3 and 2.4 give the certificate iff and exact gap; Theorem 4.3 gives an upper-value majorant; Theorem 5.1 gives the exact payoff-family cover, greatest selector, common root value and mixing closure. All original priors, Borel payoffs and boundary paths are retained.","human_contributions":"Author: Ji Ho Bae (LawAll). The previous assistant-authored restriction of the author to direction, prompting and submission is withdrawn. Detailed human-contribution wording must follow the author-approved contribution statement, not an inference from AI tool logs.","verification":"The checks documented by this Codex workflow include model-assisted mathematical/citation/exposition reviews and compilation/PDF checks. These records do not establish the absence or scope of human mathematical review. The previous assistant-generated assessment of human verification is withdrawn.","reproducibility":"Published standalone manuscript/source archive at v2.0.0 commit 31c8af23d55dc93e3f1839b2b16d34978375c4c0. Repository Prompt A applied to this new submission preparation after the original research; this chronology is disclosed. Prompt B audit before filing; public process record and scope in submission appendix. Private transcripts are not published."}
% FULL_SOLUTION_SCOPE: {"target":"For arbitrary fixed finite nonempty K,I,J, intrinsically characterize all bounded Borel payoffs admitting a value at every prior in the exact one-sided-information simultaneous-action game.","result":"Theorems 2.3, 2.4, 4.3 and 5.1; exact certificate characterization and gap, upper-value regularization and payoff-family structure.","coverage":"All finite nonempty state/action sets, all bounded Borel payoffs, every prior including boundary priors, null histories and all infinite paths; neither special-case restriction nor effective-algorithm claim.","dependencies":"TJoens-De Bock PMLR147 (2021) Eq3/Thms2,13 Borel superhedging; Durrett (2019) Thm4.6.8; elementary Fatou, compactness and conditional-expectation facts. Assumptions checked in the paper.","human_verification":"The previous assistant-generated assessment of human verification is withdrawn. Its scope must be stated by the author; AI tool records alone do not determine it.","unverified":"No identified remaining proof gap in the stated mathematical results. The pinned earlier Lean companion has its own documented scope; this submission does not extend that scope to the added v2 theorems. Catalogue acceptance remains separate.","novelty":"Primary sources and the fixed catalogue statement were checked. No exhaustive historical-priority search for the added theorems; no first-proof claim."}
% SCOPE: arbitrary fixed finite nonempty K,I,J; bounded Borel payoffs; all priors
% BASE_REVISION: 4ea46a1e7f4c898860201c820a34a34cd5a465df
\documentclass[11pt,a4paper]{article}
\usepackage[T1]{fontenc}
\usepackage{lmodern}
\usepackage[margin=25mm]{geometry}
\usepackage{amsmath,amssymb,amsthm,mathtools}
\usepackage{microtype,enumitem,xurl,hyperref}
\hypersetup{hidelinks,hypertexnames=false,
 pdftitle={Payoff families and value in games with one-sided information},
 pdfauthor={Ji Ho Bae}}
\newtheorem{theorem}{Theorem}[section]
\newtheorem{proposition}[theorem]{Proposition}
\newtheorem{lemma}[theorem]{Lemma}
\newtheorem{corollary}[theorem]{Corollary}
\theoremstyle{definition}
\newtheorem{definition}[theorem]{Definition}
\newtheorem{example}[theorem]{Example}
\theoremstyle{remark}
\newtheorem{remark}[theorem]{Remark}
\newcommand{\R}{\mathbb R}
\newcommand{\N}{\mathbb N}
\newcommand{\E}{\mathbb E}
\newcommand{\Pp}{\mathbb P}
\newcommand{\HH}{\mathcal H}
\newcommand{\co}{\operatorname{conv}}
\newcommand{\lv}{\underline v}
\newcommand{\uv}{\overline v}
\newcommand{\UU}{\mathcal U}
\newcommand{\DD}{\mathcal D}
\newcommand{\VV}{\mathcal V}
\newcommand{\ind}{\mathbf 1}
\setlength{\emergencystretch}{2em}
\allowdisplaybreaks[1]
\title{Payoff families and value in games\\with one-sided information}
\author{Ji Ho Bae\\\small LawAll}
\date{11 October 2026}
\begin{document}
\maketitle
\begin{abstract}
We study bounded Borel payoffs in repeated zero-sum games with a finite
hidden state known to one player. For fixed finite action sets, value
existence at every prior is equivalent to a condition on posterior and
capital arrays over the public history tree. The infimum of the initial capital
excess equals the minimax gap. These certificates yield an exact cover
of the value-admitting payoffs by Borel selections of pathwise closed
convex regions. Each family has a greatest payoff of Baire class at most
three, a common value at every prior, and closure under public-path
Borel mixing. Independently of value existence, every bounded Borel
payoff has a majorant of Baire class at most three with the same upper
value at every prior. The proofs use Borel superhedging duality on finite-alphabet trees
and an explicit posterior reduction. All pathwise conditions include
plays of probability zero.
\end{abstract}

\section{Model and question}\label{sec:model}

Let $K,I,J$ be nonempty finite sets with the discrete topology.
Write $\N=\{0,1,\ldots\}$ and
\[
 X=I\times J,\qquad \HH=\bigcup_{t\geq0}X^t,
 \qquad \Omega=X^{\N}.
\]
The space $\Omega$ has the product topology of the discrete topology on
$X$. A history $h\in\HH$ is public; $\varnothing$ denotes the empty
history, $hij$ appends $(i,j)$ to $h$, and $\omega|t$ is the prefix of
length $t$.

Nature draws $k\in K$ with prior $p\in\Delta(K)$ and reveals it only to
the maximizing player. At every date the players choose actions
simultaneously and then both observe the pair. Their behavioral
strategies are
\[
 \sigma:K\times\HH\longrightarrow\Delta(I),\qquad
 \tau:\HH\longrightarrow\Delta(J).
\]
The two current actions are drawn independently conditional on the
state and history. For a bounded Borel payoff $f:K\times\Omega\to\R$,
let $\gamma_{p,f}(\sigma,\tau)=\E_{p,\sigma,\tau}f(k,\omega)$ and
\[
 \lv_f(p)=\sup_\sigma\inf_\tau\gamma_{p,f}(\sigma,\tau),\qquad
 \uv_f(p)=\inf_\tau\sup_\sigma\gamma_{p,f}(\sigma,\tau).
\]
The game has a value at $p$ when these two quantities are equal;
their common value is then denoted by $v_f(p)$.

Castellon Koltun, Lehrer and Solan~\cite[Section~5]{KLS} ask for properties of
the payoff that guarantee value existence in this model. The
prior-uniform formulation in problem C73-8 of~\cite{Catalogue} asks
for a characterization of
\[
 \VV_{K,I,J}
 =\{f:K\times\Omega\to\R\text{ bounded Borel}:
           \lv_f(p)=\uv_f(p)\text{ for every }p\in\Delta(K)\}.
\]
It requires conditions on the payoff and the information structure,
beyond the equality of the defining optimizations. We give local
convex conditions and pointwise boundary inequalities, then describe
the payoff families they determine. The same arrays describe families of path-dependent payoff changes
that preserve the value at every prior.
The characterization is infinitary; no effective recognition
procedure for an arbitrary Borel payoff is asserted.

Choose $a<b$ with $a\leq f\leq b$ and put $F=(f-a)/(b-a)$. Thus
$0\leq F\leq1$. Both values transform by this positive affine map.
We write $\VV^{[0,1]}_{K,I,J}$ for the normalized class and
$\Delta_{\mathbb Q}(K)$ for the rational-coordinate priors. The
strategy spaces do not depend on $p$, including at its boundary.
Consequently,
\begin{equation}\label{eq:lipschitz}
 |\lv_F(p)-\lv_F(p')|\leq\|p-p'\|_1,\qquad
 |\uv_F(p)-\uv_F(p')|\leq\|p-p'\|_1.
\end{equation}
These bounds hold for every fixed strategy pair and survive either
order of infimum and supremum. Both values are also monotone in $F$.

\section{Local capital certificates}\label{sec:certificates}


\begin{definition}[Capital arrays at prior $p$]
At each $h\in\HH$, assign
\[
 q(h)\in\Delta(K),\quad
 A_k(h),B(h)\in[0,1]\ (k\in K).
\]
Require $q(\varnothing)=p$ and the following conditions.

\emph{Posterior children:}
\begin{equation}
 q(hij)\text{ is independent of }j;
 \quad\text{denote it by }\bar q_i(h).
 \tag{BF}\label{bf}
\end{equation}
In particular, a posterior child depends on the current informed action
but is independent of the current uninformed action, including branches
that will receive zero convex weight.

\emph{Common-column capital geometry:} define
\[
 V_j(h)=\bigl(A_k(hij)-A_k(h)\bigr)_{(k,i)\in K\times I}
 \in\R^{K\times I}.
\]
Require
\begin{equation}
 \co\{V_j(h):j\in J\}\cap\R_{-}^{K\times I}\ne\varnothing,
 \tag{UC}\label{uc}
\end{equation}
where $\R_{-}^{K\times I}$ is the coordinatewise nonpositive orthant.

\emph{Posterior and lower-capital geometry:} define
\[
 Z_i(h)=\bigl(\bar q_i(h),
        (B(hij)-B(h))_{j\in J}\bigr)\in\R^K\times\R^J.
\]
Require
\begin{equation}
 \co\{Z_i(h):i\in I\}\cap
        \bigl(\{q(h)\}\times\R_-^J\bigr)\ne\varnothing.
 \tag{LC}\label{lc}
\end{equation}

Define the bounded Borel function on \emph{every} path
\begin{equation}
 G_q(\omega)=\liminf_{t\to\infty}
       \sum_{k\in K}q_k(\omega|t)F(k,\omega).
 \tag{G}\label{g}
\end{equation}
Require the pointwise boundary conditions
\begin{align}
 \liminf_{t\to\infty}A_k(\omega|t)&\geq F(k,\omega)
       &&(k\in K,\omega\in\Omega),\tag{UB}\label{ub}\\
 \liminf_{t\to\infty}B(\omega|t)&\geq1-G_q(\omega)
       &&(\omega\in\Omega).\tag{LB}\label{lb}
\end{align}
These include paths of probability zero.
\end{definition}

\begin{definition}[Payoff certificate condition]
The payoff $F$ satisfies condition \textup{(C)} if for every
$p\in\Delta_{\mathbb Q}(K)$ and every integer $n\geq1$, there are capital
arrays as above with
\begin{equation}
 \sum_{k\in K}p_k A_k(\varnothing)+B(\varnothing)\leq1+\frac1n.
 \tag{RC}\label{rc}
\end{equation}
\end{definition}

The local conditions are finite convex-hull intersections at each
history. Their boundary inequalities compare the arrays with the payoff
on every path. We call arrays satisfying these requirements a capital
certificate; the root bound is imposed separately in (C).

\begin{theorem}[Certificate characterization]\label{main}
For arbitrary fixed nonempty finite $K,I,J$ and every bounded Borel $f$,
\[
 f\in\mathcal V_{K,I,J}\quad\Longleftrightarrow\quad
 F\text{ satisfies \textup{(C)}}.
\]
The equivalence is independent of the normalization bounds $a,b$.
\end{theorem}

We will also establish the exact local representation of the value gap.
Let $D_F(p)$ be the infimum of
$\sum_k p_k A_k(\varnothing)+B(\varnothing)-1$ over capital arrays at prior $p$
satisfying \eqref{bf}--\eqref{lb}, without \eqref{rc}. Such arrays always
exist: take $q(h)=p$ and $A_k(h)=B(h)=1$.

\begin{theorem}[Exact gap representation]\label{rep}
For every bounded Borel $F:K\times\Omega\to[0,1]$ and every prior $p$,
\[
 D_F(p)=\uv_F(p)-\lv_F(p).
\]
\end{theorem}

\section{Proof of the capital representation}\label{sec:capitalproof}
\subsection{Borel superhedging on a finite-alphabet tree}


We use equation~(3) and Theorems~2 and~13 of T'Joens and De Bock
\cite{TDB}. Equation~(3) defines the game-theoretic upper expectation
by pointwise superhedging. Theorem~2 identifies it with the axiomatic
upper expectation, and Theorem~13 identifies the latter with the
measure-theoretic upper expectation for bounded Borel payoffs.
Together they give the following form:

\begin{lemma}[Borel superhedging]\label{duality}
On a finite alphabet tree, assign each history $h$ a nonempty closed
convex set $P_h$ of next-step probability vectors. For a bounded
nonnegative Borel $W$, let $c$ be the supremum of $\E_P W$ over all laws
whose one-step kernels lie in $P_h$. For every $\delta>0$ there is a
real-valued, globally bounded-below supermartingale $M$ with $M(\varnothing)<c+\delta$,
\[
 \sup_{\pi\in P_h}\sum_{z\in X}\pi_z M(hz)\leq M(h),
 \qquad
 \liminf_t M(\omega|t)\geq W(\omega)\quad(\omega\in\Omega).
\]
\end{lemma}

The local upper expectation is $Q_h(Z)=\sup_{\pi\in P_h}\pi\cdot Z$,
so the probability sets and upper expectations agree in the terminology
of~\cite{TDB}. Both applications below use bounded Borel payoffs and
nonempty closed convex local models on the finite alphabet $X$.

\begin{lemma}[Capital can be bounded by one]\label{clip}
If $0\leq W\leq1$, the process in Lemma~\ref{duality} may be chosen in
$[0,1]$, with the same bound on its initial capital.
\end{lemma}
\begin{proof}
Any process furnished by the lemma is nonnegative at every node. Its
local inequality implies $\min_z M(hz)\leq M(h)$. If $M(h)<0$ at some
node, successively choose a child whose value is no larger than its
parent's. The resulting infinite path extending $h$ has
$\liminf M\leq M(h)<0$, contradicting its nonnegative boundary lower
bound. A finite prefix does not affect this liminf.

Now put $\widetilde M(h)=\min\{M(h),1\}$. For each $\pi\in P_h$,
concavity and monotonicity of $z\mapsto\min\{z,1\}$ give
\[
 \sum_z\pi_z\min\{M(hz),1\}
 \leq\min\left\{\sum_z\pi_zM(hz),1\right\}
 \leq\min\{M(h),1\}.
\]
Thus $\widetilde M$ satisfies the local inequality. Also
$\liminf\widetilde M=\min\{\liminf M,1\}\geq W$ on every path.
Clipping decreases initial capital, and $\widetilde M\in[0,1]$.
\end{proof}

\subsection{Posterior flow}

\begin{lemma}[Bayes-flow correspondence]\label{bayes}
Fix $p\in\Delta(K)$. Let $q:\HH\to\Delta(K)$ satisfy
$q(\varnothing)=p$ and \eqref{bf}, and let $w_i(h)\geq0$ be coefficients summing
to one with $q(h)=\sum_i w_i(h)\bar q_i(h)$. Define
\[
 x_k(i\mid h)=\frac{w_i(h)\bar q_{i,k}(h)}{q_k(h)}
 \quad\text{when }q_k(h)>0,
\]
and use any fixed action distribution when $q_k(h)=0$. This is an
informed behavioral strategy. Against every uninformed strategy $\tau$,
the public posterior at every positive-probability history is $q(h)$,
the informed action has public distribution $w(h)$, and
\begin{equation}
 \E_{p,x,\tau}F(k,\omega)=\E_{p,x,\tau}G_q(\omega).
 \tag{PI}\label{pi}
\end{equation}
Conversely every informed strategy $x$ at prior $p$ gives such arrays
with these public laws.
\end{lemma}
\begin{proof}
At a positive-$q_k(h)$ row, the barycentric identity gives
$\sum_i x_k(i\mid h)=1$. If $q_k(h)=0$, nonnegativity and that same
identity imply $w_i(h)\bar q_{i,k}(h)=0$ for every $i$. Hence
\[
 \sum_k q_k(h)x_k(i\mid h)
 =\sum_k w_i(h)\bar q_{i,k}(h)=w_i(h).
\]

Induct on the length of a positive-probability public history $h$.
At the root the posterior is $p=q(\varnothing)$. If the posterior at $h$
is $q(h)$, the next pair $(i,j)$ has conditional probability
$w_i(h)\tau_j(h)$. When this is positive, Bayes' rule gives the next
posterior
\[
 \frac{q_k(h)x_k(i\mid h)}{w_i(h)}
 =\bar q_{i,k}(h)=q_k(hij).
\]
The current factor $\tau_j(h)$ cancels because the uninformed player's
randomization does not depend on the hidden state. A zero-probability
child's assigned posterior does not affect the induction or law.
Zero-prior states are covered by the same argument.

The vector $q(\omega|t)$ is consequently the bounded posterior
martingale under every $\tau$. Its coordinates converge almost surely
to the conditional probabilities of the hidden state given the entire
public play. This is the upward conditional-expectation martingale
convergence theorem \cite[Theorem~4.6.8]{Durrett}, because the finite public-history sigma fields
generate the sigma field of the full public play. On that event,
\[
 G_q(\omega)=\sum_k q_{\infty,k}(\omega)F(k,\omega).
\]
Condition each $\mathbf1_{\{k=\ell\}}F(\ell,\omega)$ on the entire
public play, then sum over $\ell$, to obtain \eqref{pi}.

There is no exceptional-path measurability convention: at each $t$,
$\sum_k q_k(\omega|t)F(k,\omega)$ is bounded Borel, so its pointwise
liminf is bounded Borel on every path, including paths where the
posterior does not converge.

Conversely start with any informed strategy $x$ and $q(\varnothing)=p$.
At each history put
\[
 w_i(h)=\sum_k q_k(h)x_k(i\mid h).
\]
For $w_i(h)>0$ put
$\bar q_{i,k}(h)=q_k(h)x_k(i\mid h)/w_i(h)$, and for $w_i(h)=0$ put
$\bar q_i(h)=p$. Set $q(hij)=\bar q_i(h)$ for every $j$. These arrays
satisfy \eqref{bf} and the barycentric identity. At a zero-weight branch all terms
$q_k(h)x_k(i\mid h)$ are zero, so the chosen vector contributes nothing
to the barycentric identity. The preceding induction proves the claimed
public laws and \eqref{pi}.
\end{proof}

\subsection{The guarantees}

Fix capital arrays at prior $p$. Condition \eqref{uc} provides at every
$h$ coefficients $y_j(h)\geq0$, summing to one, such that
\[
 \sum_j y_j(h)A_k(hij)\leq A_k(h)
 \quad\text{for every }(k,i)\in K\times I.
\]
These coefficients define an uninformed behavioral strategy $y$; the
same distribution works for all states and all current informed
actions, as required by its information restriction.

In each state $k$, against every informed strategy $\sigma_k$, the
bounded nonnegative process $A_k(\omega|t)$ is a supermartingale.
Indeed, condition on $h$ and a pure informed action $i$, apply the
displayed inequality, and then average over that informed action.
Fatou's lemma and \eqref{ub} imply
\[
 \E_{k,\sigma_k,y}F(k,\omega)
 \leq\E\liminf_t A_k(\omega|t)
 \leq\liminf_t\E A_k(\omega|t)
 \leq A_k(\varnothing).
\]
After weighting by $p$ this proves
\begin{equation}
 \uv_F(p)\leq\sum_k p_k A_k(\varnothing).
 \label{upper}
\end{equation}

Condition \eqref{lc} supplies convex coefficients $w_i(h)\geq0$ summing
to one, simultaneously satisfying
\[
 q(h)=\sum_i w_i(h)\bar q_i(h),\qquad
 \sum_i w_i(h)B(hij)\leq B(h)\quad(j\in J).
\]
Reconstruct $x$ using these coefficients and Lemma~\ref{bayes}. Against every
$\tau$, the public next-action kernel is
$w_i(h)\tau_j(h)$. Condition \eqref{lc} makes $B(\omega|t)$ a bounded
nonnegative supermartingale: apply its inequality for each pure $j$
and average over $\tau(h)$. Fatou, \eqref{lb}, and \eqref{pi} give
\[
 \E_{p,x,\tau}(1-F)=\E(1-G_q)
 \leq\E\liminf_t B(\omega|t)\leq B(\varnothing).
\]
Consequently
\begin{equation}
 \lv_F(p)\geq1-B(\varnothing).
 \label{lower}
\end{equation}
The elementary minimax inequality, \eqref{upper}, and \eqref{lower}
yield
\begin{equation}
 0\leq\uv_F(p)-\lv_F(p)
 \leq\sum_k p_k A_k(\varnothing)+B(\varnothing)-1.
 \label{gapbound}
\end{equation}
Thus (C) gives equality at every rational prior by letting $n$ tend to
infinity.

The Lipschitz bounds~\eqref{eq:lipschitz} extend equality from the
rational priors to the entire simplex. Positive affine normalization
gives sufficiency for the original payoff.

\subsection{Sharpness of the bounds}

For arbitrary $F$ and $p$ write $L=\lv_F(p)$ and $U=\uv_F(p)$. Given
$\delta>0$, the definitions of the extrema provide behavioral strategies
$x,y$ such that
\begin{equation}
 \inf_\tau\E_{p,x,\tau}F\geq L-\delta,
 \qquad \sup_\sigma\E_{p,\sigma,y}F\leq U+\delta.
 \label{approx}
\end{equation}
No optimal strategy is assumed to exist. Obtain $q,w,\bar q$ from $x$
using Lemma~\ref{bayes}.

For each state $k$ use the local credal sets
\[
 P_h^y=\{(a_i y_j(h))_{ij}:a\in\Delta(I)\}.
\]
These are nonempty closed convex sets, being linear images of the
finite simplex. Their local upper expectation of a next-step array
$Z_{ij}$ is
\[
 \max_{i\in I}\sum_j y_j(h)Z_{ij}.
\]
Compatible precise kernels are exactly the kernels of arbitrary
state-$k$ informed behavioral responses to the fixed $y$: the marginal
on $I$ recovers $a(h)$. Therefore the compatible laws are exactly those
response laws. Let $c_k$ be their upper expectation of $F(k,\cdot)$.
Lemma~\ref{duality}, with the finite alphabet $X$ and this bounded Borel
payoff, followed by Lemma~\ref{clip}, gives $A_k\in[0,1]$ with
\[
 A_k(\varnothing)<c_k+\delta,\qquad
 \sum_j y_j(h)A_k(hij)\leq A_k(h)\quad(h\in\HH,i\in I),
\]
and with \eqref{ub} on every path. The same $y$ supplies convex
coefficients for all $k$, proving \eqref{uc} for the combined arrays.

The finitely many state components can be chosen independently, so
\begin{equation}
 \sum_k p_k c_k=\sup_\sigma\E_{p,\sigma,y}F.
 \label{statewise}
\end{equation}
For every $\sigma$, $\E_{p,\sigma,y}F\leq\sum_k p_k c_k$, so its
supremum is also at most that sum. Conversely, choose
$\varepsilon$-optimal responses separately in each positive-prior state,
combine them into one informed behavioral strategy, and let
$\varepsilon\downarrow0$. Zero-prior states have no effect. Hence
\eqref{approx} gives
\begin{equation}
 \sum_k p_k A_k(\varnothing)<U+2\delta.
 \label{acost}
\end{equation}

For the other capital process, use the public-law credal sets
\[
 P_h^w=\{(w_i(h)b_j)_{ij}:b\in\Delta(J)\}.
\]
They are again nonempty closed convex linear images of a finite
simplex. Their local upper expectation is
$\max_j\sum_i w_i(h)Z_{ij}$. Compatible precise trees correspond exactly
to uninformed strategies against $x$, because their $J$ marginal is
$b$ and their public $I$ marginal is $w$. This statement is also valid
at histories of probability zero, whose specified kernels do not
affect the root law.

The boundary payoff $W=1-G_q$ is bounded, nonnegative, and Borel.
Lemma~\ref{bayes} shows that its upper expectation over these compatible
trees equals
\[
 \sup_\tau\E(1-G_q)=1-\inf_\tau\E_{p,x,\tau}F.
\]
Apply Lemmas~\ref{duality} and \ref{clip}. They give $B\in[0,1]$ with
$\sum_i w_i(h)B(hij)\leq B(h)$ for every $j$, and satisfying
\eqref{lb}. Together with the barycentric identity, these inequalities
show that $\sum_i w_i(h)Z_i(h)$ belongs to
$\{q(h)\}\times\R_-^J$, proving \eqref{lc}. The root bound is
\begin{equation}
 B(\varnothing)<1-\inf_\tau\E_{p,x,\tau}F+\delta
 \leq1-L+2\delta.
 \label{bcost}
\end{equation}
All hypotheses of the Borel duality have thus been checked for both
applications.

Combining \eqref{acost} and \eqref{bcost} yields capital arrays with
\[
 \sum_k p_k A_k(\varnothing)+B(\varnothing)-1
 <U-L+4\delta.
\]
Taking the infimum of these costs and letting $\delta\downarrow0$, then
using \eqref{gapbound}, proves
$D_F(p)=U-L$. This establishes Theorem~\ref{rep} for every bounded
Borel payoff and every prior, without a value-existence assumption.
When the game has a value, $U=L$; choosing $4\delta<1/n$ gives
\eqref{rc} for every $n$. This proves necessity of (C), and completes
Theorem~\ref{main}. Its equivalence with value existence also proves
independence of the normalization bounds.


\section{Upper-value regularization}\label{sec:regularization}

The upper part of the proof gives a separate representation that does
not use posterior arrays. Let
\[
 \UU(F)=\{A:\HH\to[0,1]^K:\text{\eqref{uc} and \eqref{ub} hold}\}.
\]
This set has no designated prior.

\begin{proposition}\label{prop:upperdual}
For every normalized Borel payoff $F$ and every prior $p$,
\begin{equation}\label{eq:upperdual}
 \uv_F(p)=\inf_{A\in\UU(F)}p\cdot A(\varnothing).
\end{equation}
\end{proposition}
\begin{proof}
The argument for~\eqref{upper} uses only \eqref{uc} and \eqref{ub},
so it gives one inequality. The construction leading to~\eqref{acost}
provides $A\in\UU(F)$ with $p\cdot A(\varnothing)<\uv_F(p)+2\delta$
for every $\delta>0$, proving the other.
\end{proof}

We use the convention that Baire class zero consists of continuous
real-valued functions, and class at most $m+1$ consists of pointwise
limits of sequences of functions of class at most $m$.

\begin{lemma}\label{lem:baire}
For a nonempty countable family of arrays $A^d:\HH\to[0,1]^K$, the
function on $K\times\Omega$ defined by
\[
 c(k,\omega)=\inf_d\liminf_t A^d_k(\omega|t)
\]
has Baire class at most three.
\end{lemma}
\begin{proof}
For fixed $d,k,t$, the prefix function $A^d_k(\omega|t)$ is continuous.
For each $N$, its tail infimum is the pointwise limit, as $M\to\infty$,
of the continuous functions $\min_{N\leq t\leq M}A^d_k(\omega|t)$.
It has class at most one. The increasing limit over $N$ is
$\liminf_t A^d_k(\omega|t)$, of class at most two. Enumerate the
indices $d$, repeating terms if the family is finite. The infimum
over $d$ is the decreasing limit of finite minima of these class-two
functions, hence has class at most three. Finite minima preserve
each class, by applying the continuous minimum operation to the
defining sequences. The finitely many state coordinates can be
combined without raising the class.
\end{proof}

\begin{theorem}[Upper-value majorant]\label{thm:majorant}
For every Borel $F:K\times\Omega\to[0,1]$ there is a function
$C:K\times\Omega\to[0,1]$ of Baire class at most three such that
\[
 F\leq C,\qquad \uv_C(p)=\uv_F(p)\quad(p\in\Delta(K)).
\]
Every Borel $G$ with $F\leq G\leq C$ has the same upper value as $F$
at every prior.
\end{theorem}
\begin{proof}
For every $d=(r,n)$ in the countable set
\[
 \DD=\Delta_{\mathbb Q}(K)\times\N_{\geq1},
\]
choose $A^d\in\UU(F)$ with
$r\cdot A^d(\varnothing)\leq\uv_F(r)+1/n$, using~\eqref{eq:upperdual}.
Set $C(k,\omega)=\inf_d\liminf_tA^d_k(\omega|t)$. Then $F\leq C$,
Lemma~\ref{lem:baire} applies, and every $A^d$ also belongs to $\UU(C)$.
For rational $r$ and every $n$,
\[
 \uv_F(r)\leq\uv_C(r)
 \leq r\cdot A^{(r,n)}(\varnothing)\leq\uv_F(r)+1/n.
\]
Let $n\to\infty$, then use~\eqref{eq:lipschitz} to extend equality
to all priors. The assertion for $G$ follows by monotonicity.
\end{proof}

For an arbitrary bounded Borel $f$, apply the theorem to
$F=(f-a)/(b-a)$ and set $\widehat f=a+(b-a)C$. Then
$f\leq\widehat f$ and $\uv_{\widehat f}(p)=\uv_f(p)$ for every prior,
whether or not the game has a value.

\begin{corollary}\label{cor:interval}
Every $f\in\VV_{K,I,J}$ has a bounded majorant $\widehat f$ of Baire
class at most three such that every bounded Borel $g$ with
$f\leq g\leq\widehat f$ belongs to $\VV_{K,I,J}$ and satisfies
$v_g(p)=v_f(p)$ for every prior.
\end{corollary}
\begin{proof}
Normalize $f$ and apply Theorem~\ref{thm:majorant}. For $F\leq G\leq C$,
\[
 v_F(p)=\lv_F(p)\leq\lv_G(p)\leq\uv_G(p)=\uv_F(p)=v_F(p).
\]
Undo the affine normalization, setting $\widehat f=a+(b-a)C$.
\end{proof}

The majorant theorem preserves the upper value even for games without
a value. It is not a value-existence criterion. No optimality of the
bound on the Baire class is asserted.

\section{An exact cover by payoff families}\label{sec:families}

We now keep full certificates rather than only their upper arrays.
This yields families of payoffs with the same value at every prior.
Their definition uses the information and action structure, without
first choosing a payoff.

\subsection{Boundary regions}

An array family $\Theta$ assigns $(q^d,A^d,B^d)$ to every
$d=(p,n)\in\DD$. Require $q^d(\varnothing)=p$,
\eqref{bf}--\eqref{lc}, and~\eqref{rc} for each $d$.
The payoff-dependent boundary conditions have not yet been imposed.
Define
\begin{align*}
 u^d_k(\omega)&=\liminf_t A^d_k(\omega|t),&
 c_{\Theta,k}(\omega)&=\inf_{d\in\DD}u^d_k(\omega),\\
 \ell^d(\omega)&=1-\liminf_tB^d(\omega|t),&
 Q^d(\omega)&=\bigcap_{N\geq0}
       \overline{\{q^d(\omega|t):t\geq N\}}.
\end{align*}
The set $Q^d(\omega)$ is the nonempty compact cluster set in
$\Delta(K)$, defined even when the posterior sequence does not converge.
Impose the boundary compatibility condition
\begin{equation}\label{eq:bc}\tag{BC}
 \min_{r\in Q^d(\omega)}r\cdot c_\Theta(\omega)
                  \geq\ell^d(\omega)\qquad(d\in\DD,\ \omega\in\Omega).
\end{equation}
Let $\mathfrak T_{K,I,J}$ be the class of families satisfying these
requirements. All arrays take values in the same simplexes and
intervals as the capital arrays in Section~\ref{sec:certificates}.
We identify $c_\Theta$ with the payoff $(k,\omega)\mapsto c_{\Theta,k}(\omega)$.

For $\Theta\in\mathfrak T_{K,I,J}$ define
\begin{equation}\label{eq:region}
 \mathcal P_\Theta(\omega)=
 \left\{z\in[0,1]^K:
 \begin{array}{l}
 z_k\leq c_{\Theta,k}(\omega)\quad(k\in K),\\
 r\cdot z\geq\ell^d(\omega)\quad(d\in\DD,\ r\in Q^d(\omega))
 \end{array}\right\},
\end{equation}
and let $\mathcal S_\Theta$ consist of the Borel payoffs $F$ with
$F(\cdot,\omega)\in\mathcal P_\Theta(\omega)$ at every path.
The regions are intersections of closed half-spaces and the cube;
they need not be polyhedra with finitely many defining inequalities.

\begin{theorem}[Payoff-family representation]\label{thm:families}
For arbitrary fixed nonempty finite $K,I,J$,
\begin{equation}\label{eq:cover}
 \VV^{[0,1]}_{K,I,J}
       =\bigcup_{\Theta\in\mathfrak T_{K,I,J}}\mathcal S_\Theta.
\end{equation}
For each $\Theta\in\mathfrak T_{K,I,J}$ the following hold.
\begin{enumerate}[label=\textup{(\roman*)},leftmargin=*,itemsep=3pt]
\item Each $\mathcal P_\Theta(\omega)$ is nonempty, closed and convex,
and the correspondence $\omega\mapsto\mathcal P_\Theta(\omega)$
has a Borel graph. The greatest element of $\mathcal S_\Theta$
in the pointwise order is $c_\Theta$, of Baire class at most three.
\item Every $F\in\mathcal S_\Theta$ has the same value function
\begin{equation}\label{eq:familyvalue}
 v_\Theta(p)=\inf_{d\in\DD}p\cdot A^d(\varnothing)
                    \qquad(p\in\Delta(K)).
\end{equation}
In particular, $v_\Theta$ is concave and $1$-Lipschitz in the
$\ell^1$ norm.
\item If $F^1,\ldots,F^m\in\mathcal S_\Theta$ and
$\lambda:\Omega\to\Delta(\{1,\ldots,m\})$ is Borel, then
\[
 F^\lambda(k,\omega)=\sum_{a=1}^m\lambda_a(\omega)F^a(k,\omega)
             \quad\text{belongs to }\mathcal S_\Theta.
\]
Also, if $F\in\mathcal S_\Theta$ and $G$ is Borel with
$F\leq G\leq c_\Theta$, then $G\in\mathcal S_\Theta$.
\end{enumerate}
The cover in~\eqref{eq:cover} is not asserted to be disjoint.
\end{theorem}

\subsection{Proof of the representation}

For a sequence $q_t\in\Delta(K)$ with cluster set $Q$ and any fixed
$z\in\R^K$, compactness gives
\begin{equation}\label{eq:cluster}
 \liminf_t q_t\cdot z=\min_{r\in Q}r\cdot z.
\end{equation}
Every cluster point is the limit of a subsequence, so its scalar
product is at least the liminf. Conversely, a subsequence realizing
the scalar liminf has a further convergent subsequence in $\Delta(K)$,
which attains equality. Apply this identity at each path, with
$z=F(\cdot,\omega)$. The lower boundary inequalities in
\eqref{eq:region} are exactly~\eqref{lb} for all $d$, while
$F\leq c_\Theta$ is exactly~\eqref{ub} for all $d$. Thus, for Borel $F:K\times\Omega\to[0,1]$,
\begin{equation}\label{eq:selectionequiv}
 F\in\mathcal S_\Theta
 \quad\Longleftrightarrow\quad
 (q^d,A^d,B^d)\text{ certifies }F\text{ at its prior for every }d.
\end{equation}

Condition~\eqref{eq:bc} says that $c_\Theta(\omega)$ itself belongs
to $\mathcal P_\Theta(\omega)$. Conversely, any point $z$ of that
region implies~\eqref{eq:bc}, since $z\leq c_\Theta$ and every
$r\in\Delta(K)$ has nonnegative coordinates. This proves
nonemptiness; closedness and convexity follow from~\eqref{eq:region}.
Lemma~\ref{lem:baire} shows that $c_\Theta$ is Borel and has the
claimed Baire class. It is therefore the greatest member of
$\mathcal S_\Theta$.

The graph is Borel as well. By~\eqref{eq:cluster}, its defining
conditions are the coordinate upper bounds and the countably many
inequalities
\[
 \liminf_t q^d(\omega|t)\cdot z\geq\ell^d(\omega).
\]
For each $d,t$, the function $(\omega,z)\mapsto q^d(\omega|t)\cdot z$
is continuous. The displayed comparisons are consequently Borel;
no measurable-selection theorem is needed.

If $F\in\mathcal S_\Theta$, equivalence~\eqref{eq:selectionequiv}
and the root bounds give (C), hence value existence by
Theorem~\ref{main}. Conversely, if $F\in\VV^{[0,1]}_{K,I,J}$,
choose a certificate satisfying the root bound for each $d=(p,n)$.
The upper boundaries imply $F\leq c_\Theta$. For every cluster point
$r\in Q^d(\omega)$, the lower boundaries and~\eqref{eq:cluster} give
\[
 r\cdot c_\Theta(\omega)\geq r\cdot F(\cdot,\omega)
                                      \geq\ell^d(\omega).
\]
Hence (BC) holds, $\Theta\in\mathfrak T_{K,I,J}$, and
$F\in\mathcal S_\Theta$. This proves~\eqref{eq:cover}.

To identify the value, fix $F\in\mathcal S_\Theta$. Each $A^d$
belongs to $\UU(F)$, and so gives the upper bound
$v_F(p)\leq p\cdot A^d(\varnothing)$ at \emph{every} prior, regardless
of the prior assigned to $q^d$. For rational $p$, use $d=(p,n)$,
\eqref{rc} and~\eqref{lower} to obtain
\[
 p\cdot A^{(p,n)}(\varnothing)
 \leq1-B^{(p,n)}(\varnothing)+1/n\leq v_F(p)+1/n.
\]
This proves~\eqref{eq:familyvalue} at rational priors. The right-hand
side is concave and $1$-Lipschitz, being the infimum of linear
functions whose coefficient vectors lie in $[0,1]^K$.
Equation~\eqref{eq:lipschitz} and density extend equality to all priors.
The expression is independent of $F$, proving the common-value claim.

Finally, pathwise convexity proves the mixing assertion. The weights
depend on the public path, with the same weights for all states.
If $F\leq G\leq c_\Theta$, nonnegativity of $r$ preserves every
lower inequality in~\eqref{eq:region}; the upper inequalities are
already assumed. This proves the last assertion and completes the
proof of Theorem~\ref{thm:families}.

\section{Strategies and payoff compatibility}\label{sec:examples}

The local arrays retain the game's strategic constraints. Their
weights reconstruct behavioral strategies as in Section~\ref{sec:capitalproof}.
It is useful to make this relation explicit before distinguishing
the payoff families from ordinary value level sets.

\subsection{Fixed-strategy costs}

Fix $p,F$ and an uninformed strategy $y$. Let $\mathcal A_y(F)$ be
the upper arrays satisfying \eqref{ub} and
\[
 \sum_jy_j(h)A_k(hij)\leq A_k(h)\quad(h\in\HH,\ k\in K,\ i\in I).
\]
For an informed strategy $x$, take the recursive posterior extension
$(q^x,w^x)$ in Lemma~\ref{bayes}, and let $\mathcal B_x(F)$ be
the $[0,1]$-valued arrays satisfying~\eqref{lb} for $q^x$ and
$\sum_iw_i^x(h)B(hij)\leq B(h)$ for every $h,j$.
The two applications of superhedging in Section~\ref{sec:capitalproof},
now with these strategies fixed, give
\begin{align}
 \inf_{A\in\mathcal A_y(F)}p\cdot A(\varnothing)
       &=\sup_\sigma\gamma_{p,F}(\sigma,y),\label{eq:fixedupper}\\
 \inf_{B\in\mathcal B_x(F)}B(\varnothing)
       &=1-\inf_\tau\gamma_{p,F}(x,\tau).\label{eq:fixedlower}
\end{align}
The lower bounds follow from the corresponding supermartingale
guarantees; Lemma~\ref{duality} gives arbitrarily close upper
bounds. Thus eliminating the capital variables recovers
\[
 D_F(p)=\inf_{x,y}\left\{
       \sup_\sigma\gamma_{p,F}(\sigma,y)
       -\inf_\tau\gamma_{p,F}(x,\tau)\right\}.
\]
This is an identity of infima, with no claim of attainment. The
representation does not remove strategy search. Its further content
is the description of the whole payoff family allowed by common
pathwise bounds, including its greatest member and value-preserving
changes.

\subsection{Common optimal strategies do not determine a family}

\begin{example}\label{ex:matching}
Let $K$ be a singleton, $I=J=\{0,1\}$, and
\[
 F(\omega)=\ind_{\{i_0=j_0\}},\qquad G(\omega)=1-F(\omega).
\]
Both games have value $1/2$, and their optimal strategies are exactly
those with a uniform first action; later actions are immaterial.
Nevertheless, no $\mathcal S_\Theta$ contains both payoffs.
\end{example}

\begin{proof}
If both belonged to one family, every array in that family would
satisfy the boundaries for both. Since $q\equiv1$, they would imply
\[
 \liminf_tA(\omega|t)\geq\max(F,G)=1,\qquad
 \liminf_tB(\omega|t)\geq\max(1-F,1-G)=1.
\]
The supermartingale guarantees force $A(\varnothing)=B(\varnothing)=1$.
The certificate cost is then $1$, contradicting the root bound for
$n=2$.

Individually, each payoff has a certificate of cost zero. For $F$,
put $A(\varnothing)=B(\varnothing)=1/2$, and after the first pair
keep $A=F$ and $B=1-F$. Uniform local weights satisfy the two convex
conditions at the root; at later histories the arrays are constant.
Interchanging $F$ and $G$ gives the other certificate.
\end{proof}

The same distinction appears in the path-dependent mixture with
weights $\lambda_1=F$ and $\lambda_2=1-F$: it equals $1$ everywhere.
Its value is $1$, whereas mixing within one $\mathcal S_\Theta$
would preserve $1/2$. The families therefore record compatibility
of pathwise payoff bounds beyond equality of values and optimal
strategies.

\subsection{A fixed-strategy infimum need not be attained}

\begin{example}\label{ex:nonattainment}
Let $K,I$ be singletons, let $J=\{0,1\}$, and fix the uninformed
strategy $y$ to choose each action with probability $1/2$ at every
history. Let $F$ indicate the path consisting entirely of zeros.
The infimum in~\eqref{eq:fixedupper} is zero and is not attained.
\end{example}

\begin{proof}
Identify $I\times J$ with $J$ and write $hj$ for the child reached
by action $j$. The payoff has expectation zero under the fixed
fair-coin law. If a
nonnegative upper array had root zero, the inequalities
$(A(h0)+A(h1))/2\leq A(h)$ would force every node to have value zero.
This violates the boundary on the all-zero path.

For each $N\geq1$, start instead with $A(\varnothing)=2^{-N}$.
Before time $N$, double the capital after a zero and set it to zero
after a one. After time $N$, keep it constant. This $[0,1]$-valued
array has the required boundary and satisfies the local fair-coin
inequality with equality. Its root values tend to zero.
\end{proof}

The nonattainment in this example concerns the \emph{fixed} strategy
$y$. In the unrestricted certificate problem, cost zero is attained:
take $A(\varnothing)=0$, put $A(h)=1$ at nonempty all-zero histories
and $A(h)=0$ elsewhere, and set $B\equiv1$. The upper weights always
choose action $1$. This satisfies every certificate condition.

Theorem~\ref{thm:families} concerns the greatest payoff within a
\emph{specified} family, not a greatest payoff among all games with
the same value. It also supplies neither an optimal strategy nor a
zero-cost certificate. The approximation index is retained in the
definition of $\Theta$. Within those limits, the local geometry and
boundary compatibility give an exact description of the full
prior-uniform class and of the payoff changes preserved by each
family.

\begin{thebibliography}{9}
\bibitem{KLS}
Gil Bar Castellon Koltun, Ehud Lehrer and Eilon Solan.
\emph{The Maxmin Value of Repeated Games with Incomplete Information
on One Side and Tail-Measurable Payoffs}.
\href{https://arxiv.org/pdf/2512.01581v1}{arXiv:2512.01581v1} (2025).
Definition~2.1, pp.~6--7; Section~5, p.~18.

\bibitem{Catalogue}
\emph{Payoff functionals guaranteeing a value under one-sided
information}, problem C73-8, Economics catalogue in
\texttt{mehmetmars7/Erdosproblems-llm-hunter}.
\href{https://github.com/mehmetmars7/Erdosproblems-llm-hunter/blob/b1f633121074e16579fa7fd05314e5f504d29bae/attacks/open_problems/economics/statements/C73-8.tex}{Fixed statement at commit \texttt{b1f633121074}}.

\bibitem{TDB}
Natan T'Joens and Jasper De Bock.
\emph{Global Upper Expectations for Discrete-Time Stochastic
Processes: In Practice, They Are All The Same!}
\href{https://proceedings.mlr.press/v147/t-joens21a/t-joens21a.pdf}{Proceedings of Machine Learning Research \textbf{147}, 310--319 (2021)}.
Equation~(3), p.~314; Theorems~2 and~13, pp.~315 and~318.
Extended version:
\href{https://arxiv.org/pdf/2102.13075v2}{arXiv:2102.13075v2},
Appendix~A.2, p.~14 (proof of Theorem~13).

\bibitem{Durrett}
Rick Durrett. \emph{Probability: Theory and Examples}, fifth edition
(2019). \href{https://sites.math.duke.edu/~rtd/PTE/PTE5_011119.pdf}{Author's version of 11 January 2019},
Theorem~4.6.8, p.~247.
\end{thebibliography}

\section*{Submission history, scope and verification}

\paragraph{Original target and result.}
This is a substantive revision of the hand-proof attempt in
\href{https://github.com/mehmetmars7/Erdosproblems-llm-hunter/pull/134}{PR~134}.
Theorem~\ref{main} retains every finite nonempty state and action set,
every bounded Borel payoff and every prior in C73-8. Theorem~\ref{rep}
identifies the exact gap. Theorem~\ref{thm:families} adds the payoff-family
cover, greatest payoff, common value and path-dependent mixing;
Theorem~\ref{thm:majorant} adds upper-value regularization without a
value-existence assumption. These structural conclusions support the
claim that the characterization supplies payoff information beyond
restating the two optimization problems. No effective decision procedure
or removal of all strategy search is claimed. The original manuscript,
source statement and historical submissions are preserved.

\paragraph{External dependencies and context.}
The exact dependencies are stated and applied in Section~\ref{sec:capitalproof}:
T'Joens--De Bock's Borel superhedging duality and Durrett's upward
conditional-expectation theorem, together with the elementary probability
and compactness arguments given in the proofs. The cited primary sources
and fixed catalogue statement were consulted directly. Their roles and
hypotheses are identified in the paper. This revision makes no exhaustive
literature-priority or first-proof claim.

\paragraph{Author and AI tools.}
The author is Ji Ho Bae (LawAll). The previous assistant-written
restriction of the author's role, and its assessment of human
mathematical review, are withdrawn. AI tool logs are not a complete
account of the collaborative research process. Detailed contribution
wording is to follow the author's own contribution statement.

The documented tools include OpenAI Codex and ChatGPT Pro, as well as
Anthropic Claude. This revision's Codex workflow used
\texttt{gpt-6-astra} with \texttt{ultra} reasoning and separate
\texttt{high}-reasoning review tasks. The underlying model version of
the supplied ChatGPT Pro response was not exposed.
The \href{https://github.com/jbaelaw/borel-value-certificates-lean-20261011/blob/adfa513ba33427af32e3bca4a01f0a86a1086db4/README.md}{Lean companion's attribution}
credits Claude Opus 5.5 and Claude Fable 5.1 (Claude Code) with the Lean
development. Its documented scope includes the original certificate
characterization, exact gap, Borel superhedging duality, convex-hull
equivalence and bounded-payoff normalization. This contribution is
part of the project's history and is not attributed to Codex.

\paragraph{Checks and remaining verification.}
The complete mathematical text underwent separate model-assisted
mathematical, citation and exposition reviews. The reviews checked the
original quantifiers, zero-prior states, null paths, the cluster-set
identity, Baire bounds, the root-value formula and both examples.
They identified no remaining proof gap in the stated results. The
Bayes-flow root assumption, fixed-strategy nonattainment scope,
normalization argument and theorem hierarchy were made explicit.
The desktop compiler and Tectonic passed, and all eleven pages of the
published manuscript were visually inspected. These statements describe the checks performed by the reported tools;
they do not determine the scope of the author's mathematical review.
\href{https://github.com/mehmetmars7/Erdosproblems-llm-hunter/pull/155}{PR~155}
records the Claude-developed Lean companion. The scope of that pinned
companion and the added v2 results are distinguished in this submission.

\paragraph{Process and reproducibility.}
The mathematical research and revision preceded this repository-prompt
application. Prompt~A from the repository's contribution prompts was
applied prospectively on 11 October 2026 to preparation of this new
catalogue submission, using the original statement, completed manuscript,
contribution policy and actual provenance. This does not retroactively
establish that Prompt~A preceded the original research. Prompt~B was
applied to the completed submission before filing. It found no new
mathematical gap. Its contribution-label recommendation was subsequently
corrected because it did not establish the full collaborative history. The original Prompt~A timing difference
is disclosed for maintainer assessment; no retrospective compliance
with the original A-before-research sequence is claimed. Private conversations and credentials are not
included in the public artifacts.

The published mathematical text is reproduced here without changing its
proofs; this appendix supplies the catalogue's declarations.
The pinned \href{https://github.com/jbaelaw/borel-game-value-characterization/tree/31c8af23d55dc93e3f1839b2b16d34978375c4c0}{source commit}
contains the standalone \href{https://github.com/jbaelaw/borel-game-value-characterization/blob/31c8af23d55dc93e3f1839b2b16d34978375c4c0/paper/Proof.tex}{TeX},
\href{https://github.com/jbaelaw/borel-game-value-characterization/blob/31c8af23d55dc93e3f1839b2b16d34978375c4c0/paper/Proof.pdf}{PDF}
and review records. The
\href{https://github.com/jbaelaw/borel-game-value-characterization/releases/tag/v2.0.0}{v2.0.0 release}
and \href{https://doi.org/10.5281/zenodo.23290245}{Zenodo version DOI}
archive this revision. The
\href{https://doi.org/10.5281/zenodo.23286393}{previous Zenodo version}
retains the earlier materials. The v1 tag, assets and both prior PRs
are not replaced or automatically closed.

\section{Completion Estimate}
\noindent\textbf{COMPLETION: 100\%}
Full-solution claim for the original C73-8 scope: the certificate iff
and payoff-family structure cover all of its stated parameters and
quantifiers. No unresolved mathematical step is identified in the
present proof. This estimate concerns claimed coverage, not proof
confidence or catalogue acceptance.

\section*{Final status}
\textbf{SOLVED (LLM claim).} Submitted for review as a new substantive
version. Source-catalogue status and community-review verdicts are
unchanged.

\end{document}
